\section*{Choice assumptions and the original ultraproduct source}
\addcontentsline{toc}{section}{Choice assumptions and the original ultraproduct source}

\begin{proofboundary}{def:p03-ultrafilter-choice-boundary}{Boolean prime ideals, ultrafilters, and Choice}
Over \(\mathsf{ZF}\), the Boolean Prime Ideal Theorem is equivalent to the
Ultrafilter Lemma and is strictly weaker than the full Axiom of Choice, relative
to the usual consistency assumptions \cite{HowardRubin1998}.  Applied to the
cofinite filter on \(\mathbb N\), it suffices for the existence of a
nonprincipal ultrafilter.

The exact inherited development is written under the handbook's ambient
\(\mathsf{AC}\) convention.  Its assumption map is:
\begin{itemize}
  \item extending the cofinite filter to a nonprincipal ultrafilter needs only
        the Ultrafilter Lemma;
  \item the inherited type-space compactness proof invokes the already stated
        first-order compactness theorem in its ambient metatheory rather than
        claiming a new \(\mathsf{ZF}+\mathsf{BPI}\) calibration;
  \item the existential step in Los's theorem selects witnesses across the
        relevant index set, so the displayed proof retains full Choice;
  \item compactness via ultraproducts uses an ultrafilter on finite subtheories
        and also selects one model for each finite subtheory, so the displayed
        proof likewise retains full Choice; and
  \item the Keisler--Shelah theorem remains cited under the ambient Choice
        convention and is not weakened here.
\end{itemize}
The distinctions are metatheoretic assumption boundaries, not new axioms of
the object theories.
\index{Boolean Prime Ideal Theorem}\index{Ultrafilter Lemma}\index{Axiom of Choice}
\end{proofboundary}

\begin{proofboundary}{def:p03-los-source-correction}{Source correction for Los's theorem}
The exact inherited sentence cites the retained bibliography key
\texttt{Los1955}, whose work is \emph{On the Extending of Models I}.  That is a
different paper.  The original ultraproduct construction and fundamental
theorem occur in Los's \emph{Quelques remarques, th\'eor\`emes et probl\`emes
sur les classes d\'efinissables d'alg\`ebres}, pages 98--113 of
\emph{Mathematical Interpretation of Formal Systems}
\cite{Los1955Ultraproducts}.  The retained key remains exact bibliography
evidence; this correction supplies the controlling source for the theorem.
\end{proofboundary}
